\documentclass[11pt, a4paper]{article}

% ── Packages ──────────────────────────────────────────────────────
\usepackage[utf8]{inputenc}
\usepackage{amsmath, amssymb}
\usepackage{graphicx}
\usepackage{hyperref}
\usepackage{booktabs}
\usepackage{natbib}
\usepackage[margin=2.5cm]{geometry}

% ── Title ─────────────────────────────────────────────────────────
\title{Quantum Error Mitigation Benchmarking Suite:\\
       ZNE, PEC, and CDR on Real IBM Hardware}
\author{John George Alexander}
\date{\today}

\begin{document}
\maketitle

\begin{abstract}
% TODO: Write after completing all experiments
% Should summarise: (1) what was benchmarked, (2) key findings,
% (3) when each method helps/fails, (4) practical recommendations.
\end{abstract}

% ── 1. Introduction ──────────────────────────────────────────────
\section{Introduction}
% Background on NISQ noise, why error mitigation matters
% Brief overview of ZNE, PEC, CDR

% ── 2. Methods ───────────────────────────────────────────────────
\section{Methods}

\subsection{Zero-Noise Extrapolation (ZNE)}
% Gate folding, Richardson extrapolation

\subsection{Probabilistic Error Cancellation (PEC)}
% Quasi-probability decomposition, sampling overhead

\subsection{Clifford Data Regression (CDR)}
% Near-Clifford training, linear regression correction

% ── 3. Experimental Setup ────────────────────────────────────────
\section{Experimental Setup}
% IBM backend details, circuit descriptions, depth range, repetitions

% ── 4. Results ───────────────────────────────────────────────────
\section{Results}
% Noise-vs-depth curves, comparison tables, confidence intervals

% ── 5. Discussion ────────────────────────────────────────────────
\section{Discussion}
% When does each method win? Failure cases. Overhead trade-offs.

% ── 6. Conclusion ───────────────────────────────────────────────
\section{Conclusion}
% Summary, practical recommendations, future work

% ── References ───────────────────────────────────────────────────
\bibliographystyle{unsrt}
\bibliography{references}

\end{document}
